\subsubsection{Machine du modèle}
\label{subsubsec:machine}

La machine définit l'état dynamique du système, les invariants de sécurité et les événements possibles :

\begin{lstlisting}[style=EventB, caption=Extrait de la machine Event-B]
MACHINE TrafficLightControl

SETS
    DIRECTIONS = {NORTH, SOUTH, EAST, WEST};
    COLORS = {RED, YELLOW, GREEN};
    PEDESTRIAN_SIGNALS = {WALK, DONT_WALK};
    PHASES = {PHASE_NS, PHASE_EW, TRANSITION_NS_TO_EW, TRANSITION_EW_TO_NS}

VARIABLES
    traffic_lights,      // État des feux pour les véhicules
    pedestrian_signals,  // État des feux pour les piétons
    current_phase,       // Phase actuelle du cycle
    timer,               // Temps écoulé dans la phase actuelle
    queue_length,        // Longueur des files d'attente
    waiting_time,        // Temps d'attente cumulé
    total_waiting_time   // Temps d'attente total (à minimiser)

INVARIANTS
    // Typage des variables
    traffic_lights ∈ DIRECTIONS → COLORS &
    pedestrian_signals ∈ DIRECTIONS → PEDESTRIAN_SIGNALS &
    current_phase ∈ PHASES &
    timer ∈ ℕ &
    queue_length ∈ DIRECTIONS → ℕ &
    waiting_time ∈ DIRECTIONS → ℕ &
    total_waiting_time ∈ ℕ &
    
    // Propriété de sécurité fondamentale:
    // Les feux opposés ne peuvent pas être verts simultanément
    ∀d1, d2 · (d1 ∈ DIRECTIONS & d2 ∈ DIRECTIONS & 
              ((d1 = NORTH & d2 = EAST) or (d1 = NORTH & d2 = WEST) or
               (d1 = SOUTH & d2 = EAST) or (d1 = SOUTH & d2 = WEST))) ⇒
              ¬(traffic_lights(d1) = GREEN & traffic_lights(d2) = GREEN)
END
\end{lstlisting}

\subsection{Propriétés de sécurité}
\label{subsec:securite}

Le modèle Event-B garantit deux propriétés de sécurité essentielles :

\begin{enumerate}
    \item \textbf{Absence de conflits} : Les feux de directions perpendiculaires ne peuvent jamais être verts simultanément, évitant ainsi les collisions.
    \item \textbf{Sécurité des piétons} : Les piétons ne peuvent traverser que lorsque les véhicules correspondants sont arrêtés (feux rouges).
\end{enumerate}

Ces propriétés sont formellement vérifiables grâce aux invariants définis dans la machine Event-B.

\subsection{Événements du système}
\label{subsec:evenements}

Le modèle Event-B définit plusieurs événements qui modélisent la dynamique du système :

\begin{itemize}
    \item \texttt{VehicleArrives} : Arrivée d'un nouveau véhicule dans une file
    \item \texttt{Tick} : Avancement du temps et mise à jour des temps d'attente
    \item \texttt{SwitchFromNS\_to\_Transition}, \texttt{SwitchFromTransition\_to\_EW}, etc. : Changements de phases des feux
    \item \texttt{VehiclesDeparture} : Départ des véhicules aux feux verts
    \item \texttt{OptimizePhases} : Calcul des durées optimales des phases
\end{itemize}

\begin{figure}[h]
\centering
\begin{tikzpicture}[node distance=2cm, auto]
    % Définir les styles
    \tikzstyle{phase} = [circle, draw, fill=blue!20, text width=2.5cm, text centered, minimum height=1cm]
    \tikzstyle{transition} = [circle, draw, fill=yellow!40, text width=2.5cm, text centered, minimum height=1cm]
    \tikzstyle{line} = [draw, -latex', thick]
    
    % Phases
    \node [phase] (ns) {Phase Nord-Sud (Vert)};
    \node [transition, right=of ns] (ns_to_ew) {Transition NS→EW (Jaune)};
    \node [phase, below=of ns_to_ew] (ew) {Phase Est-Ouest (Vert)};
    \node [transition, left=of ew] (ew_to_ns) {Transition EW→NS (Jaune)};
    
    % Connections
    \path [line] (ns) -- (ns_to_ew);
    \path [line] (ns_to_ew) -- (ew);
    \path [line] (ew) -- (ew_to_ns);
    \path [line] (ew_to_ns) -- (ns);
    
    % Durées optimisées
    \node [draw, dashed, fit=(ns) (ns_to_ew) (ew) (ew_to_ns), inner sep=1cm, label=above:{\textbf{Cycle avec durées optimisées}}] {};
\end{tikzpicture}
\caption{Cycle des phases des feux tricolores}
\label{fig:phases}
\end{figure}

\section{Modèle mathématique pour l'optimisation des temps d'attente}
\label{sec:optimisation}

Pour minimiser le temps d'attente des véhicules, nous avons implémenté un modèle mathématique adaptatif qui ajuste les durées des phases en fonction des longueurs des files d'attente.

\subsection{Algorithme d'optimisation adaptatif}
\label{subsec:algorithme}

L'algorithme que nous avons implémenté est basé sur un modèle adaptatif qui calcule les durées optimales des phases proportionnellement aux longueurs des files :

\begin{algorithm}
\caption{Optimisation adaptative des durées de phases}
\begin{algorithmic}[1]
\State $q_{NS} \gets$ longueur de la file Nord-Sud
\State $q_{EW} \gets$ longueur de la file Est-Ouest
\State $q_{total} \gets q_{NS} + q_{EW}$
\If{$q_{total} > 0$}
    \State $d_{NS} \gets d_{min} + (d_{max} - d_{min}) \times \frac{q_{NS}}{q_{total}}$
    \State $d_{EW} \gets d_{min} + (d_{max} - d_{min}) \times \frac{q_{EW}}{q_{total}}$
\Else
    \State $d_{NS} \gets d_{default}$
    \State $d_{EW} \gets d_{default}$
\EndIf
\end{algorithmic}
\end{algorithm}

où:
\begin{itemize}
    \item $d_{min}$ : Durée minimale d'une phase (10 secondes)
    \item $d_{max}$ : Durée maximale d'une phase (60 secondes)
    \item $d_{default}$ : Durée par défaut (12 secondes)
\end{itemize}
